\documentclass[11pt]{article}

\usepackage[margin=1in]{geometry}
\usepackage{graphicx}
\usepackage{booktabs}
\usepackage{amsmath}
\usepackage{hyperref}
\usepackage{xcolor}
\usepackage{caption}
\usepackage{subcaption}

\graphicspath{{figures/}}

\title{Understanding Persona Learning via Probe Monitors}
\author{Sriram and Eric}
\date{}

\begin{document}
\maketitle

% ============================================================================
\section{Introduction}
% ============================================================================

% Can we predict the behavior of a fine-tuned model before it is visible
% from its generated output?


% ============================================================================
\section{Method}
% ============================================================================

% Probe extraction pipeline, fingerprinting formula, scoring conventions.

% Key formula: fingerprint = lora(own_responses) - clean(own_responses)
% Each model generates its own text AND scores through its own model.

\subsection{Trait Vector Extraction}

\subsection{Fingerprinting}

% The fingerprint delta for variant $v$ on trait $t$:
% \begin{equation}
%     \Delta_t^v = \text{cos}(\bar{a}_t^v, \mathbf{w}_t) - \text{cos}(\bar{a}_t^{\text{clean}}, \mathbf{w}_t)
% \end{equation}
% where $\bar{a}_t^v$ is the mean response activation at trait $t$'s layer,
% and $\mathbf{w}_t$ is the probe vector.

\subsection{Probe Validation}

Probe scores are compared against an independent LLM judge (GPT-4o-mini, scoring 0--100) on responses from 7 LoRA-finetuned tonal personas (Qwen3-4B). The probe measures activation-space cosine deltas on clean text prefilled through each LoRA; the judge scores the LoRA's own generated text. Agreement between these fundamentally different signals---one internal, one behavioral---validates the probes as tracking real behavioral change.

\begin{table}[h]
\centering
\begin{tabular}{lrrrrr}
\toprule
Trait & Spearman $\rho$ & p-value & Pearson $r$ & Probe AUROC & Judge AUROC \\
\midrule
nervous\_register      & +0.647 & 3.6e-84 & +0.709 & 1.000 & 1.000 \\
bureaucratic           & +0.606 & 2.0e-71 & +0.953 & 1.000 & 1.000 \\
angry\_register        & +0.594 & 6.8e-68 & +0.822 & 1.000 & 0.985 \\
disappointed\_register & +0.581 & 2.0e-64 & +0.523 & 0.985 & 0.999 \\
confused\_processing   & +0.571 & 8.0e-62 & +0.624 & 0.997 & 1.000 \\
mocking                & +0.275 & 1.3e-13 & +0.262 & 0.764 & 0.985 \\
\bottomrule
\end{tabular}
\caption{Per-trait correlation between probe cosine deltas and LLM judge scores (all 700 responses pooled). AUROC measures discrimination of matching vs.\ non-matching persona LoRAs.}
\label{tab:probe-validation}
\end{table}


% ============================================================================
\section{Emergent Misalignment Detection}
% ============================================================================

\subsection{Variant Fingerprints}

% Top trait deltas for Turner et al. EM variants.

% \begin{figure}[h]
% \centering
% \includegraphics[width=\textwidth]{em_grouped_bars.png}
% \caption{Top 23 trait deltas by variant (best steering layers).}
% \label{fig:em-bars}
% \end{figure}

\subsection{Fingerprint Correlation}

% Variant-to-variant similarity across top traits.

% \begin{figure}[h]
% \centering
% \includegraphics[width=0.7\textwidth]{em_correlation_heatmap.png}
% \caption{Pearson $r$ between variant fingerprints.}
% \label{fig:em-correlation}
% \end{figure}

\subsection{Early Detection via Checkpoints}

% Cosine similarity to final fingerprint over training.

% \begin{figure}[h]
% \centering
% \includegraphics[width=0.8\textwidth]{em_early_detection.png}
% \caption{Fingerprint cosine to final checkpoint during training.}
% \label{fig:em-early}
% \end{figure}


% ============================================================================
\section{Reward Hacking Detection}
% ============================================================================

\subsection{Population-Level Trait Shifts}

% Cohen's d: RH vs baseline.

% \begin{figure}[h]
% \centering
% \includegraphics[width=0.8\textwidth]{rh_cohens_d.png}
% \caption{Top 10 / Bottom 10 traits by Cohen's $d$ (RH vs Baseline).}
% \label{fig:rh-cohens-d}
% \end{figure}

\subsection{Cross-Seed Consistency}

% \begin{figure}[h]
% \centering
% \includegraphics[width=\textwidth]{rh_3seed.png}
% \caption{Trait shift comparison across 3 seeds.}
% \label{fig:rh-seeds}
% \end{figure}

\subsection{Onset Dynamics}

% Per-token trait trajectories centered on hack onset.

% \begin{figure}[h]
% \centering
% \includegraphics[width=\textwidth]{rh_onset_dynamics.png}
% \caption{Top traits around reward hack onset. Red = rising, blue = dropping.}
% \label{fig:rh-onset}
% \end{figure}

% \begin{figure}[h]
% \centering
% \includegraphics[width=\textwidth]{rh_onset_heatmap.png}
% \caption{All traits around hack onset (model-specific signal).}
% \label{fig:rh-onset-heatmap}
% \end{figure}


% ============================================================================
\section{Discussion}
% ============================================================================


\end{document}
